%%%%%%%%%%%%%%%%%%%%%%%%%%%%%%%%%%%%%%%%%%%%%%%%%%%%%
%			HLAVIČKA								%
%%%%%%%%%%%%%%%%%%%%%%%%%%%%%%%%%%%%%%%%%%%%%%%%%%%%%
\documentclass[openany,12pt]{memoir}

\usepackage[utf8]{inputenc} 
\usepackage[czech]{babel}
\usepackage[T1]{fontenc}
\usepackage[top=1.5cm, bottom=2cm, left=2cm, right=2cm]{geometry}  % --> NASTAVENÍ OKRAJŮ
\usepackage{fancyhdr}
\usepackage{graphicx}
%\usepackage{lmodern}
\usepackage{tgschola}  % --- Silnější font
\usepackage{xwatermark}
\usepackage{xcolor}
\usepackage{changepage}
\usepackage{newunicodechar}



%%%%%% Package na zpěvník
\usepackage[full]{leadsheets}%http://mirrors.nic.cz/tex-archive/macros/latex/contrib/leadsheets/leadsheets_en.pdf   --> dokumentace	
\definesongtitletemplate{empty}{} 
\setchords{
format = \bfseries,   %tučné akordy
minor = {mi},% 
input-notation = {german},%
output-notation = {german}%
}
\definesongtitletemplate{empty}{} 

\newlength{\drop}
\newwatermark[allpages,color=red!50,angle=0,scale=2, xpos=0,ypos=0]{\includegraphics[width=5cm]{obr/pozadi2.jpg}} %--> dvojka na pozadí


%%%% Vlastní příkazy
\newcounter{Slokočet}   %Automatické číslování slok
\newcommand{\mezera}{\vspace*{0.5cm}}   %Horizontální odsazení slok
\newcommand{\stred}{5.2cm}   %%% Na zarovnání slok doprostřed, pozn. automatičtější zarovnávání na střed nejde
\newcommand{\refren}{\mezera \noindent \textbf{R:} } %refrén
\newcommand{\sloka}{\mezera \noindent \addtocounter{Slokočet}{1} \arabic{Slokočet}. } 	%sloka, která se automaticky čísluje
\newcommand{\ssloka}{\mezera \noindent}  % vlastní číslo sloky
\newcommand{\carka}{,\:}
\newcommand{\predehra}{\mezera \hspace*{-1.85cm} \textbf{Předehra:} }
\newcommand{\m}[1]{\color{white}{#1}}  %Pro akordy
\newcommand{\ap}{'} %Pro apostrof
\newcommand{\elipsa}{\kern\fontdimen3\font} %Příkaz pro lepší zacházení s výpustkami (=...); je to vpodstatě jen mezera mezi tečkama výpustky

%%%%%%%%%%%%%%%%%%%%%%%%%%%%%%%%%%%%%%%%%%%%%%%%%%%%%%
%			NÁVOD									 %
%%%%%%%%%%%%%%%%%%%%%%%%%%%%%%%%%%%%%%%%%%%%%%%%%%%%%%
%1. Věci v hlavičce IGNOROVAT
%2. Píseň psát do prostoru mezi \begin{song} a \end{song}
%3. další řádek se značí dvěma odsazeníma (= dvakrát stisknout enter)
%4. \refren vždy na začátku refrenu a \sloka na začátek sloky (automaticky se čísluje)
%  \ = alt gr + q ; [/] = alt gr f/g ; {/} = alt gr + b/n; ^ = alt gr + 3 + mezera
%Cokoliv napíšete do ^{  } se bude brát jako akord
%když se toto bude dotýkat nějakého slova (nebude mezi tím a slovem mezera)
%tak se akordy zjeví nad slovem, ale pište to před slovo
%Když se to nedotýká slova, tak akord lítá ve vzduchu a vytiskne se větší mezera
%První možnost je asi preferovanější
%5. Akordy stačí psát jen do první sloky, když se nezmění -- kytaristi to zvládnou



%\usepackage{subfiles}
%</++>

\begin{document}

\pagestyle{empty}
%\addtocontents{toc}{\protect\thispagestyle{empty}}
%\tableofcontents \thispagestyle{empty}\newpage
\newgeometry{top=1.5cm, bottom = 0cm, left = 2cm, right = 2cm}

\newpage

%Pro kompilaci je důležitý, 	aby ten input a \end{document} zůstali na posledních
%dvou řádcích!!!!
\begin{song}{title=\predtitle\centering Lano co k nebi nás poutá \\\large Traband  \vspace*{-0.3cm}}  %% sem se napíše jméno songu a autor
\begin{centerjustified}
\nejvetsi

\sloka
Já ^{F\z }sedával v přístavu, ^*{Gmi}popíj el kořalu, ^{F}s holkama ^{\z C7}laškoval.

A ^{F}bylo mi fuk, co je, ^{Gmi}hlavně když fajfka mi ^{F\z C7}doutná.

Co ^{F\,\,}bylo už není, ^*{Gmi}všech no mý jmění jsem ^{A7\,\,}dávno ^{\z Dmi}rozfofroval.

Jsme ^{F\,\,}silný, jak silný je ^{Gmi}lano, co k nebi nás ^{F C7 F}poutá.

\sloka
Ale najednou zmatek, když vešel ten chlápek, na mou duši! 

Objedná si drink a sedne si vedle do kouta.

Pak se nakloní ke mně a povídá jemně: \uv{Matouši!

Jsme silný, jak silný je lano, co k nebi nás poutá.}

\sloka
Já povídám: \uv{Pane, odkud se známe? Esli se nemýlíte?

A co je vám do mě, starýho mrchožrouta?}

On na to: \uv{Pojď, dej se na moji loď, má jméno \emph{Eternité}.} 

Jsme silný, jak silný je lano, co k nebi nás poutá.\\


^{Dmi D7}

\sloka
^{G}Ty jeho slova se ^{Ami\,\,\,}zařízly do mě ^{G\,\,\,\,}jako bys břitvou ^{D7\,\,\,\,\,\,}šmik,

^{G\,\,\,\,}jako když po ránu ^{Ami\,\,\,\,}vzbudí tě křik ^{G\,\,\,\,\, D7\,\,}kohouta.

^{G}Tak povídám: ,,Jdem!“ A ^{Ami}ještě ten den stal se ^{H7}ze mě ^{\,\,\,\, Emi}námořník. 

^{G}Jsme silný, jak silný je ^{Ami}lano, co k nebi nás ^{G D7 G}poutá.

\sloka
Tak zvedněme kotvy a napněme plachty, vítr začíná vát! 

Černý myšlenky vymeťme někam do kouta.

Hudba ať hraje o dobytí ráje, teď není čeho se bát.

/: Jsme silný, jak silný je lano, co k nebi nás poutá. :/\\

\end{centerjustified}
\setcounter{Slokočet}{0}
\end{song}

\end{document}
